\subsection*{29. Points of Parity Are the Ticket to Enter the Arena}

When a buyer evaluates an unfamiliar alternative, there is a minimal set of non-negotiable expectations. A cloud tool without adequate privacy controls, a training program with unreliable teaching, or a physical product that lacks basic safety cannot rescue itself through 1 extraordinary differentiator.

Positioning should communicate both \emph{parity} and \emph{difference} in the correct order. First, establish enough category membership to be evaluated; then reveal the distinctive value. A founder who aggressively distances the company from its category may force buyers to ask whether fundamental standards have been met. Conversely, a brand that lists only category basics can be mistaken for a generic substitute.

A practical worksheet has 3 columns: \textbf{must be credible at}, \textbf{must win on}, and \textbf{can deliberately sacrifice}. Use the customer's actual criteria rather than the founder's favorite features.

\subsection*{30. Which Market Category Should We Claim?}

Categories reduce the buyer's cognitive burden by providing expectations and reference prices. Calling something \emph{enterprise accounting software} signals different capabilities from \emph{a spreadsheet add-on}. Calling something \emph{a university degree} evokes specific legal and accreditation requirements; an unaccredited startup must not appropriate that claim.

Four common framing choices present themselves:

\begin{enumerate}
    \item \textbf{Compete inside an established category:} Easier recognition and budget allocation; direct head-to-head comparison.
    \item \textbf{Specify a subcategory:} Preserve familiarity while narrowing around a job: for example, a production AI design-review studio rather than a comprehensive AI bootcamp.
    \item \textbf{Use an adjacent category:} Change the buyer's reference set, such as specialized technical consulting versus generic education, provided the service truly fits.
    \item \textbf{Attempt a novel category:} Potential to influence the standards by which buyers compare, balanced by the substantial cost of explaining the need, language, budget, and adoption path.
\end{enumerate}

New founders frequently overestimate the fourth option. \textbf{Inventing a phrase is not creating a category.} A market category requires recognition beyond the company: buyers use it, analysts and partners understand it, comparable products emerge, budgets form, and standards become shared. Scholarly work on market category formation distinguishes emergent categories from categories deliberately promoted by actors [R6].

\subsection*{31. The Category Comprehension Tax}

Suppose a founder announces: ``We are creating the world's first autonomous cognitive capability ecosystem.'' The sentence sounds visionary, but a buyer cannot easily infer what to purchase, whose budget it uses, what problem it solves, or which alternatives it replaces. The startup must now educate customers about the problem, category, firm, and purchasing procedure before obtaining a sale.

That learning burden is an economic cost. If educating each buyer requires 6 calls, technical demos, and a complex procurement exercise, the category may be too expensive for a company with minimal capital to establish. Start with recognizable language tied to current spending and reveal the novel system through delivery.

Category creation becomes plausible when a discontinuity (new technology, regulation, workflow, social norm, or infrastructure) makes a distinctly new result attainable, and the startup can demonstrate it before the market has settled on a label. Even then, prospective category leadership is an option, not a right earned through self-description.

\subsection*{32. Category Design and the Danger of Category Captivity}

An established category makes an offering easier to understand, but it also defines what buyers assume they are buying. An education company positioned primarily as \emph{placement training} attracts customers whose fundamental purchase criterion is near-term employment. Later insisting that the same business is a research-oriented academic institution is difficult if customer expectations, faculty, pricing, and marketing have been optimized for placement promises.

Conversely, a technically rigorous new institution can be too abstract if it describes itself only in terms of scientific civilization-building when the early buyer needs a specific practical capability. The solution is a \textbf{two-layer positioning architecture}: a clear, constrained commercial proposition at the product level and an enduring, honest institutional identity above it. Neither should promise assets that do not exist.

The category chosen for the first product should be judged on 3 criteria: does it win the current comparison, is it truthful about what exists, and does it avoid unnecessarily constraining the institution's future credible evolution?
